# cockpit_glow.tex - which parts of the cabin are lit from the inside
# (Phase 25 stage C), and the first texture in this game that is not a colour.
#
# ⚑ IT IS A MASK, SO ITS COORDINATES ARE NOT ITS OWN. Every rect below is
# copied from cockpit.tex's "Instrument glow" op, and it has to stay copied:
# this file says WHICH pixels of that texture glow, in that texture's uv space,
# so a strip moved there and not here stops glowing and a strip moved here and
# not there lights up bare panelling. They are two halves of one decision and
# the only thing keeping them together is this paragraph.
#
# ⚑ WHY IT EXISTS AT ALL, WHICH IS A FIX AND NOT A LOOK. models.toml lifted the
# whole cockpit with a flat `emissive = 0.10`, for a stated and correct reason -
# vacuum ambient is 1.2%, which is right for a hull seen from outside and pitch
# black for a room the player sits in. But a flat lift is the wrong shape for
# what cockpit.tex says it is drawing: "a cabin is dark, with the light coming
# from the instruments". A scalar cannot say "these five strips and nothing
# else"; a second texture can. The `emissive` floor is KEPT beside it, so this
# adds the instruments rather than redistributing the cabin.
#
# ⚑ Red channel only is what the shader reads, but this is written as white
# because the cooker's encoder is BC1 RGB - a mask stored in one channel would
# still cost three, and grey values here would read as partial glow, which is a
# feature nothing wants yet.
size = [256, 256]

[[op]]
kind = "fill"
color = [0, 0, 0]

# The five instrument strips, exactly as cockpit.tex paints them in amber.
[[op]]
kind = "rects"
color = [255, 255, 255]
rects = [
  [148, 197, 26, 4],
  [173, 66, 26, 4],
  [105, 194, 26, 4],
  [40, 105, 26, 4],
  [50, 156, 26, 4],
]
